% Generated by build.mjs from ledger.mjs. Do not edit by hand.
\begin{interpretation}{I-1}{Non-triviality}
\item[Issue] S1 requires a \enquote{non-trivial} theory but does not define the term.
\item[Working reading] $\mathcal T_G$ is not isomorphic to (i) the one-dimensional theory $\mathcal H=\mathbb C\Omega$ (already excluded by YM-14), (ii) the Wick-polynomial theory of $\dim\mathfrak g$ free massless vector fields (already excluded by YM-13, since it is gapless), or (iii) a generalized free (Gaussian) theory. Recommended certificate: prove YM-11 for at least one connected $n$-point function, $n\ge 3$, whose leading perturbative coefficient is non-zero.
\item[Positive track] The certificate settles YM-04 and part of YM-11 together.
\item[Negative track] A result of the form \enquote{only trivial theories exist} must exclude (iii) explicitly.
\end{interpretation}

\begin{interpretation}{I-2}{Compact simple gauge group}
\item[Issue] As an abstract group $SU(N)$ is not simple (centre $\mathbb Z_N$). The physics usage is: compact connected Lie group with simple, non-abelian Lie algebra. Different global forms (e.g.\ $SU(N)$ and $SU(N)/\mathbb Z_N$) share the local operator content on $\mathbb R^4$ and differ in line operators and $\theta$-periodicity \cite{AST13}.
\item[Working reading] The quantifier ranges over all compact simple Lie algebras $\mathfrak g\in\{A_n(n\ge1),B_n(n\ge2),C_n(n\ge3),D_n(n\ge4),E_6,E_7,E_8,F_4,G_2\}$, each with any compact connected $G$ such that $\operatorname{Lie}(G)=\mathfrak g$.
\item[Positive track] A method that works only for $SU(3)$, or only for $SU(N)$, does not suffice: all four classical families and all five exceptional algebras are required.
\item[Negative track] A counterexample for one single $\mathfrak g$ suffices.
\end{interpretation}

\begin{interpretation}{I-3}{Theta parameter}
\item[Issue] The Lagrangian (1) of S1 has no topological $\theta$-term, although the quantum theory admits one.
\item[Working reading] The target is the $\theta=0$ theory.
\item[Positive track] $\theta=0$ suffices (existential quantifier).
\item[Negative track] Unless CMI's reading restricts to $\theta=0$, a non-existence result would also have to cover $\theta\neq0$. Flag for expert confirmation.
\end{interpretation}

\begin{interpretation}{I-4}{Short-distance agreement}
\item[Issue] \enquote{Agree at short distances with the predictions of asymptotic freedom and perturbative renormalization theory, as described in textbooks} does not specify which correlation functions or to which order.
\item[Working reading] For gauge-invariant local fields, Schwinger functions at non-coinciding points, scaled towards a point, have asymptotic expansions in the running coupling $\bar g(\mu)$ that match renormalized perturbation theory order by order. Minimal certificate: leading-logarithmic agreement including the universal coefficients $b_0,b_1$ (\idref{A-009}).
\item[Positive track] The order of agreement must be stated explicitly in the main theorem.
\item[Negative track] Ambiguity here is a candidate for S2 \S5(c)(ii) (reformulation), not for the Prize.
\end{interpretation}

\begin{interpretation}{I-5}{Axiom set}
\item[Issue] \enquote{At least as strong as those cited in [45, 35]}: [45] gives the Wightman axioms, [35] is OS~I \emph{and} OS~II. The reconstruction claim of OS~I is false without the linear growth condition (\idref{A-002}).
\item[Working reading] Wightman axioms as in Definition~\ref{def:wightman}; on the Euclidean route the OS axioms including the linear growth condition of \cite{OS75}. A Haag--Kastler net alone is not sufficient, because the requirement cites [45, 35], not [24].
\item[Positive track] Prove OS0$'$ (linear growth) explicitly; it is the axiom most often forgotten.
\item[Negative track] \enquote{The axioms are too restrictive} is an opinion, not a counterexample (Section~\ref{sec:negative}).
\end{interpretation}

\begin{interpretation}{I-6}{Form of publication}
\item[Issue] S2 speaks of \enquote{the Proposed Solution} published in a Qualifying Outlet and does not say how a series of papers is treated.
\item[Working reading] A series of refereed papers is acceptable in principle, but one main paper must contain the complete logical chain (the \enquote{Clay corollary}) with exact references to the published components.
\item[Positive track] No essential step may live only in code, data or supplements (YM-22).
\item[Negative track] Same rule.
\end{interpretation}
